\documentclass[11pt]{article}
\pdfoutput=1
\usepackage{amsmath}
\usepackage{hyperref}
\usepackage{geometry}
\usepackage{listings}
\usepackage[dvipsnames]{xcolor}
\usepackage{booktabs}
\usepackage{longtable}
\usepackage{array}
\usepackage{parskip}
\usepackage{fancyhdr}
\usepackage{titlesec}
\usepackage{mdframed}
 
\geometry{
  letterpaper,
  top=25mm, bottom=25mm,
  left=25mm, right=25mm
}
 
% ---------------------------------------------------------------
% Code listing style
% ---------------------------------------------------------------
\definecolor{codebg}{rgb}{0.95,0.95,0.95}
\definecolor{codeframe}{rgb}{0.7,0.7,0.7}
\definecolor{commentcolor}{rgb}{0.4,0.4,0.4}
\definecolor{keywordcolor}{rgb}{0.1,0.3,0.7}
 
\lstset{
  backgroundcolor=\color{codebg},
  basicstyle=\small\ttfamily,
  breaklines=true,
  frame=single,
  rulecolor=\color{codeframe},
  tabsize=4,
  columns=flexible,
  keepspaces=true,
  showstringspaces=false,
  commentstyle=\color{commentcolor}\itshape,
  keywordstyle=\color{keywordcolor},
  numbers=none,
  xleftmargin=6pt,
  xrightmargin=6pt,
  aboveskip=8pt,
  belowskip=8pt
}
 
\newcommand{\code}[1]{\texttt{#1}}
 
% ---------------------------------------------------------------
% Note / Warning boxes
% ---------------------------------------------------------------
\newmdenv[
  linewidth=1pt,
  leftmargin=0pt, rightmargin=0pt,
  innerleftmargin=8pt, innerrightmargin=8pt,
  innertopmargin=6pt, innerbottommargin=6pt,
  skipabove=8pt, skipbelow=8pt
]{notebox}
 
\newmdenv[
  backgroundcolor=red!8,
  linewidth=1pt,
  leftmargin=0pt, rightmargin=0pt,
  innerleftmargin=8pt, innerrightmargin=8pt,
  innertopmargin=6pt, innerbottommargin=6pt,
  skipabove=8pt, skipbelow=8pt
]{warnbox}
 
% ---------------------------------------------------------------
% Headers / footers
% ---------------------------------------------------------------
\pagestyle{fancy}
\fancyhf{}
\fancyhead[L]{\small Magic2 Architecture Notes}
\fancyhead[R]{\small v1.0 -- DRAFT}
\fancyfoot[C]{\small\thepage}
\renewcommand{\headrulewidth}{0.4pt}
 
% ---------------------------------------------------------------
% Section formatting
% ---------------------------------------------------------------
\titleformat{\section}{\large\bfseries}{{\thesection}}{0.8em}{}[\vspace{2pt}\hrule\vspace{4pt}]
\titleformat{\subsection}{\normalsize\bfseries}{{\thesubsection}}{0.8em}{}
\titleformat{\subsubsection}{\normalsize\itshape\bfseries}{{\thesubsubsection}}{0.8em}{}
 
% ---------------------------------------------------------------
\begin{document}
 
% ---------------------------------------------------------------
% Title page
% ---------------------------------------------------------------
\begin{titlepage}
\vspace*{20mm}
{\Huge \textbf{Magic2}}\\[4mm]
{\Large Architecture and Developer Notes}\\[2mm]
{\large Internal Document --- Not for End Users}
\vspace{8mm}
\hrule
\vspace{8mm}
{\large Jean~Thierry-Mieg, Danielle~Thierry-Mieg, Greg~Boratyn}\\[2mm]
National Library of Medicine, National Institutes of Health\\
8600 Rockville Pike, Bethesda MD 20894, USA\\[2mm]
\href{mailto:mieg@ncbi.nlm.nih.gov}{\texttt{mieg@ncbi.nlm.nih.gov}}
\vspace{8mm}
 
\noindent Source code: \href{https://github.com/NLM-DIR/Magic2}{\texttt{https://github.com/NLM-DIR/Magic2}}
 
\vfill
\small National Library of Medicine $\cdot$ National Institutes of Health\\
\textit{This document describes internal implementation details.
For user-facing documentation see the Magic2 User Guide.}
\end{titlepage}
 
\newpage
\tableofcontents
\newpage
 
% ===============================================================
\section{Purpose of This Document}
% ===============================================================
 
This document is intended for developers working on the Magic2 source
code, not for end users.  It describes the internal pipeline
architecture, the parallelism model, the NUMA-aware process model,
the GPU code path, and other implementation details that are not
relevant to running alignments but are essential for understanding,
debugging, or extending the code.
 
For user-facing documentation, see the Magic2 User Guide
(\texttt{wsa\_doc/magic2\_userguide.pdf}).
 
% ===============================================================
\section{Pipeline Architecture}
% ===============================================================
 
\subsection{Overview}
 
Magic2 is structured as a deep pipeline of independent processing
layers.  Data flows through the pipeline as discrete blocks; all
layers execute concurrently and self-balance their workloads.  This
design is inspired by the channel and goroutine model of the Go
programming language, implemented in C using a custom channel library
(see Section~\ref{sec:channels}).
 
The key architectural properties are:
 
\begin{itemize}
  \item All pipeline layers run simultaneously.  There is no master
        thread dispatching work; each layer pulls from its input
        channel and pushes to its output channel independently.
  \item Work is divided into \emph{data blocks} of at most
        \code{bMax} megabases each.  The number of blocks
        simultaneously in flight is \code{nBlocks}.
  \item Each layer contains \code{nAgents} parallel agents.
        All agents within a layer are interchangeable and process
        whichever block arrives next.
  \item RAM consumption is controlled by a \emph{dropper} mechanism
        at the head of the pipeline that withholds new blocks when
        memory pressure is high, rather than by pre-allocating a
        fixed buffer.
  \item Input size is unbounded: blocks flow in and out continuously,
        so a run of 100~Gbases is handled identically to a run of
        1~Gbase.
\end{itemize}
 
\subsection{Pipeline layers}
 
The main alignment pipeline consists of the following layers, in order:
 
\begin{center}
\begin{tabular}{clp{7.5cm}}
\toprule
\textbf{Layer} & \textbf{Name} & \textbf{Function} \\
\midrule
1 & Reader      & Reads and decompresses input files or SRA streams;
                  emits read blocks. \\
2 & Seeder      & Extracts seeds from reads; looks up seeds in the
                  index (CPU or GPU). \\
3 & Sorter      & Radix-sorts seed hits to group candidates by target
                  position. \\
4 & Extender    & Extends seed hits to full-length alignments;
                  resolves introns on the CPU. \\
5 & Filter      & Applies score, length, and error-rate filters;
                  resolves multi-mappers. \\
6 & Writer      & Formats and writes BAM, \code{.hits}, wiggles,
                  intron tables, and count tables. \\
\bottomrule
\end{tabular}
\end{center}
 
\begin{notebox}
\textbf{Note:} The GPU accelerates layers 2 and 3 (seeding and
sorting).  Layer 4 (extension), which requires semantically complex
splice-junction logic, remains on the CPU in the current
implementation.  See Section~\ref{sec:gpu}.
\end{notebox}
 
\subsection{Key internal parameters}
 
\begin{center}
\begin{tabular}{lllp{6cm}}
\toprule
\textbf{Variable} & \textbf{Flag} & \textbf{Default} & \textbf{Meaning} \\
\midrule
\code{nAgents}     & \code{--nA} & $3/2 \times n_\text{CPU}$
                   & Agents per pipeline layer. \\
\code{nBlocks}     & \code{--nB} & $2 \times n_\text{CPU}$
                   & Simultaneous blocks in flight. \\
\code{bMax}        & \code{--bMax} & 3~MB
                   & Maximum megabases per block. \\
\code{max\_threads}& \code{--max\_threads} & hardware-derived
                   & Hard cap on UNIX threads. \\
\code{tStep}       & \code{--step} & 6 (large genomes)
                   & Seed step used during index construction. \\
\code{rStep}       & \code{--rStep} & \code{tStep}/2
                   & Seed step used during read alignment. \\
\code{NN}          & \code{--NN} & 16
                   & Number of seed file partitions. \\
\bottomrule
\end{tabular}
\end{center}
 
\begin{warnbox}
\textbf{Warning:} \code{max\_threads} must not be set below 64.
The channel architecture requires a minimum level of concurrency;
an artificially low thread cap will cause the pipeline to stall
waiting for virtual threads that cannot be scheduled.
\end{warnbox}
 
% ===============================================================
\section{The Channel Library}
\label{sec:channels}
% ===============================================================
 
\subsection{Motivation}
 
Standard C parallelism (pthreads, OpenMP) is based on shared memory
and explicit locking.  Magic2 instead uses a message-passing model
inspired by Go channels: data blocks are passed between layers through
typed channels; no layer directly reads or writes another layer's
memory.
 
This design was chosen because:
\begin{itemize}
  \item It eliminates entire classes of race conditions: a block is
        owned by exactly one layer at a time.
  \item It makes the pipeline self-documenting: the channel topology
        is the architecture diagram.
  \item It scales naturally: adding agents to a layer requires no
        changes to the surrounding code.
\end{itemize}
 
\subsection{Channel semantics}
 
Each channel is a bounded FIFO queue.  A layer blocks on a read if
the channel is empty and blocks on a write if the channel is full.
The dropper at the head of the pipeline monitors RAM usage and delays
injection of new blocks when memory is constrained, providing
back-pressure throughout the system without any layer needing to be
aware of memory state.
 
\subsection{Virtual threads}
 
Each agent and each channel endpoint runs in its own virtual thread.
Virtual threads are lightweight: the system may have thousands of them
scheduled over a smaller number of UNIX threads.  The \code{max\_threads}
parameter caps the number of UNIX threads, not the number of virtual
threads.
 
% ===============================================================
\section{NUMA-Aware Process Model}
\label{sec:numa}
% ===============================================================
 
\subsection{NUMA detection}
 
On startup, Magic2 counts the number of NUMA nodes by scanning
\code{/sys/devices/system/node}.  If the directory is absent
(non-NUMA kernel, container environment) or contains only one node,
the count is treated as 1 and no re-spawning occurs.
 
\begin{lstlisting}[language=C]
static int count_numa_nodes (void)
{
    DIR *d = opendir ("/sys/devices/system/node") ;
    if (!d) return 1 ;
    int count = 0, dummy ;
    struct dirent *e ;
    while ((e = readdir (d)))
        if (sscanf (e->d_name, "node%d", &dummy) == 1)
            count++ ;
    closedir (d) ;
    return count > 0 ? count : 1 ;
}
\end{lstlisting}
 
\subsection{Re-spawn mechanism}
 
When more than one NUMA node is detected, Magic2 identifies the least
loaded node and re-spawns itself using \code{numactl --cpunodebind=N
--membind=N}, where \code{N} is the chosen node.  This ensures that
all memory allocations for the pipeline are local to the CPUs that
will process them, avoiding the latency penalty of cross-node memory
access.
 
The re-spawned child receives the flag \code{--numactl} in its
argument list.  This flag signals that re-spawning has already
occurred and suppresses a second re-spawn, preventing infinite
recursion.
 
\subsection{Implications for debugging}
\label{sec:debug}
 
When re-spawning occurs, a debugger attached to the parent process
will not follow execution into the child.  Breakpoints set before
the \code{exec} call will never be hit.
 
Two command-line flags are provided to work around this:
 
\begin{center}
\begin{tabular}{lp{10cm}}
\toprule
\textbf{Flag} & \textbf{Effect} \\
\midrule
\code{--debug}
  & Suppresses re-spawning and forces \code{--nAgents 1 --nBlocks 1},
    giving a fully serial execution path.  Verbose diagnostic messages
    are enabled automatically.  Use this for step-by-step debugging. \\
\code{--numactl}
  & Suppresses re-spawning only.  Full parallelism is retained.
    Use this when profiling with \code{perf} or \code{vtune}, or when
    you need breakpoints but also need realistic concurrency. \\
\bottomrule
\end{tabular}
\end{center}
 
\textbf{Typical gdb sessions:}
\begin{lstlisting}
# Fully serial, easiest to step through:
gdb --args magic2 --debug -x IDX -i test.fa -o /tmp/dbg/
 
# Breakpoints work, full parallelism retained:
gdb --args magic2 --numactl -x IDX -i test.fa -o /tmp/dbg/
\end{lstlisting}
 
\begin{notebox}
\textbf{Note:} On single-node machines, in containers, and on
non-NUMA kernels, re-spawning never occurs and breakpoints work
without any special flag.
\end{notebox}
 
% ===============================================================
\section{GPU Acceleration}
\label{sec:gpu}
% ===============================================================
 
\subsection{Scope}
 
The GPU accelerates the two most computationally intensive layers of
the pipeline: seed extraction and lookup (layer~2) and radix sorting
of seed hits (layer~3).  These operations are embarrassingly parallel
and map well onto GPU execution models.
 
Layer~4 (alignment extension) remains on the CPU.  Extension requires
splice-junction logic, gap penalties, and intron model evaluation that
are semantically complex and do not parallelise efficiently on a GPU
in their current form.  This may change in future versions.
 
\subsection{Internal binary datasets}
 
Several index data structures have GPU-specific variants that differ
in memory layout from the CPU versions.  These are generated at index
construction time when a GPU compiler is available.  The flag
\code{--noJump}, which controls whether jumper entries are inserted
into the index, is a compile-time internal switch relevant only to the
GPU build; it is not exposed in the user-facing or developer help.
 
\subsection{Fallback}
 
If no CUDA-capable GPU is present at runtime, Magic2 falls back
silently to the CPU path for all layers.  No flag is needed.
 
% ===============================================================
\section{Seed Parameters}
% ===============================================================
 
\subsection{Index construction step (\code{tStep})}
 
During index construction, a seed is extracted from the target every
\code{tStep} bases.  \code{tStep} is the \emph{maximum} distance
between consecutive seeds; because the best seed among the next
\code{tStep} words is selected, the average distance is
\code{tStep}/2.
 
The default is 1 for targets smaller than 1~Mb and 6 for larger
genomes.  Values of 8--10 reduce index RAM and construction time at
a modest sensitivity cost; values of 1--4 increase sensitivity at
higher RAM and CPU cost.
 
\subsection{Alignment step (\code{rStep})}
 
During read alignment, seeds are extracted from the query every
\code{rStep} bases.  By default \code{rStep = tStep / 2}.
 
\code{rStep} is exposed in \code{--devhelp} as \code{--rStep} for
benchmarking purposes.  It should not be changed in production use.
 
\subsection{Seed file partitioning (\code{NN})}
 
The seed index is split into \code{NN} partitions to allow parallel
lookup.  The allowed values are powers of 2 from 1 to 64.
When \code{seedLength} exceeds 16, the minimum required value of
\code{NN} is $4^{(\text{seedLength} - 16)}$; this minimum is always
enforced automatically.
 
% ===============================================================
\section{The TSF Format}
% ===============================================================
 
\subsection{Overview}
 
TSF (Tag-Sample Format) is the internal tabular format used by Magic2
for all multi-sample output: gene expression counts, intron support
tables, poly-A addition sites, QC metrics, and coverage summaries.
 
\subsection{Line structure}
 
Each line contains four tab-separated fields:
 
\begin{center}
\begin{tabular}{clp{8cm}}
\toprule
\textbf{Column} & \textbf{Content} & \textbf{Notes} \\
\midrule
1 & Run name  & Identifies the sample or run. \\
2 & Tag       & Feature identifier (gene name, intron coordinates, etc.). \\
3 & Format    & Type code string describing the data columns
                (e.g.\ \code{i} = integer, \code{f} = float,
                \code{t} = text). \\
4+ & Data     & One or more data values matching the format code. \\
\bottomrule
\end{tabular}
\end{center}
 
Lines need not be sorted.  TSF files from different runs can be
concatenated directly; the TSF library merges them by joining on
column 1, column 2, or both.
 
\subsection{Type codes}
 
The format field (column 3) is a string of single-character type
codes, one per data column:
 
\begin{center}
\begin{tabular}{cl}
\toprule
\textbf{Code} & \textbf{Type} \\
\midrule
\code{i} & Integer \\
\code{f} & Float \\
\code{t} & Text \\
\bottomrule
\end{tabular}
\end{center}
 
For example, a line with format \code{iif} has two integer columns
followed by one float column.
 
\subsection{Merging and aggregation}
 
The TSF library provides a merge operation that joins lines from
multiple files on a specified key column and aggregates the data
columns.  Available aggregation operators include sum, min, max, and
others.  The merged output can be emitted as a TSF file or as a
two-dimensional table (rows = tags, columns = runs) suitable for
import into R or Python.
 
The user-facing merge command is:
\begin{lstlisting}
magic2 --mergeTSF run1/genes.tsf run2/genes.tsf -o all_genes.txt
\end{lstlisting}
 
\subsection{Relation to user guide}
 
The TSF format is mentioned briefly in the Magic2 User Guide in the
context of output files.  The internal structure, type codes, and
library API are documented here only, as they are relevant to
developers who wish to parse TSF output programmatically or extend
the TSF library.
 
% ===============================================================
\section{Known Limitations and Future Work}
% ===============================================================
 
\begin{itemize}
  \item GPU acceleration of the extension layer (layer~4) is not yet
        implemented.  The splice-junction and gap model would need to
        be restructured for efficient GPU execution.
  \item The \code{--strand} override flag is specified in the
        documentation and help text but requires implementation in
        the alignment strategy selection code.
  \item The \code{-i} flag currently accepts comma-separated file
        lists; the documented behaviour (single file or paired pair
        only, comma lists refused with a clear error message) requires
        a code change to enforce.
\end{itemize}
 
% ===============================================================
\section*{Acknowledgements}
% ===============================================================
 
We thank David Landsman and Tom Madden for their support.
This work was supported by the Intramural Research Program of the
National Library of Medicine, National Institutes of Health.
 
\end{document}
 
